
% \begin{frame}[fragile]
% 	\frametitle{Conclusion}

% 	\begin{tabular}{|c|c|c|c|}
% 		\diagbox[width=142pt, height=41pt,font=\small, innerrightsep=1em]{Contributions}{Limitations} & foo & bar & baz \\
% 		G-cartes                                                                                      &     &     &     \\
% 		Règles Jerboa                                                                                 &     &     &     \\
% 		Détection évènements                                                                          &     &     &
% 	\end{tabular}
% \end{frame}


\begin{frame}
	\frametitle{Contributions}
	Trois contributions concrétisent notre démarche~:
	\begin{itemize}
		\item la formalisation des changements topologiques
		\item la mise en œuvre d'un mécanisme de réévaluation à base de règles avec
		      \begin{itemize}
			      \item la définition d'un mécanisme de nomination persistante à deux
			            couches~;
			      \item l'appariement des entités sur la base de leurs évolutions et en
			            tenant compte des modifications~;
			      \item la définition de stratégie pour configurer le processus de rejeu
		      \end{itemize}
		\item une extension de ce mécanisme pour le rejeu de processus comprenant des scripts
	\end{itemize}
\end{frame}

\begin{frame}
	\frametitle{Perspectives}
	\begin{itemize}
		\item<1-> Compléter l'édition avec le déplacement d'opération~;
		\item<2-> Intégration des scripts au mécanisme de rejeu~;
		\item<3-> Développement de scripts de rejeu pour la modélisation en archéologie~;
		\item<4> Annotation temporelle des opérations~;
\item
	\end{itemize}
\end{frame}

%%% Local Variables:
%%% mode: LaTeX
%%% TeX-master: "../../main"
%%% End:
